%! Random Sampling Methods — Unit 3, Lesson 3.2 — Student Packet Cover Sheet
\documentclass[12pt]{article}
\usepackage{apstats-article}
\usepackage{apstats-boxes}
\usepackage{ltablex}
\keepXColumns

\begin{document}

% Full-bleed navy banner. \coverbanner MEASURES the title block and sizes the
% band to it, so a lesson title long enough to wrap grows the band rather than
% running out the bottom of it.
\coverbanner{Unit 3}{Lesson 3.2 \quad Random Sampling Methods}

% NAMESTRIP: this is the ONLY component that carries the name row.
\namedateperiod
\vspace{0.18in}

\boxguard[14]
\begin{learningtargetbox}
\small
% GRADUAL RELEASE: the vocabulary is taught in the guided notes, so the targets name it
% plainly. The AP Learning Objectives (DAT-2.C, DAT-2.D) stay in the teacher-facing plan.
\begin{enumerate}[leftmargin=1.5em, itemsep=5pt]
  \item \ldots{} describe how to choose a \textbf{simple random sample} with a random number
        generator or a table of random digits, and tell sampling \textbf{with replacement}
        from sampling \textbf{without replacement}.
  \item \ldots{} identify a \textbf{stratified random sample} and a \textbf{cluster sample}
        from the description of a study.
  \item \ldots{} explain, in context, why one sampling method is a better choice than another.
\end{enumerate}
\end{learningtargetbox}

\vspace{0.14in}

\boxguard[14]
\begin{tocbox}
\small
% A CONTENTS LIST ONLY: no Score column and no Total row. Whether a component counts is said in
% its description: AP Practice is extra credit, scored 1 to 5; the Homework is scored.
\renewcommand{\arraystretch}{1.5}
\begin{tabularx}{\linewidth}{@{} c l X @{}}
\rowcolor{sky}
\textbf{\#} & \textbf{Component} & \textbf{Description} \\
\midrule
1 & Warm-Up & Lakeview High's survey of 80 students: names in a hat, seniors, and advisory
                 groups \\
2 & Guided Notes \& Practice & Simple random, stratified, and cluster samples at Lakeview, and
                 which plan to choose; then parking at an apartment complex \\
3 & AP Practice & Farms, textbook spending, and an airplane survey: AP-style multiple choice and
                 one free-response set --- extra credit, scored 1 to 5 \\
4 & Homework & A hospital's nurses: name each plan, describe a simple random sample, and choose
                 between strata and clusters --- \textbf{scored 1 to 5, due the first class after
                 two study halls} \\
\end{tabularx}
\end{tocbox}

\vspace{0.14in}

\boxguard[8]
\begin{remindbox}
\small
A \textbf{simple random sample} gives every group of the right size the same chance --- but just
by chance it can come back lopsided. \textbf{Strata} are groups that are alike inside, like the
grades at a school: sample from \emph{every} one, so no group is left out. \textbf{Clusters} are
groups that each look like the whole population, like mixed advisory groups: take \emph{all} of
\emph{some}, which is cheaper. Before you choose, ask: are these groups alike inside, or does
each one look like the whole?
\end{remindbox}

\end{document}
